
\documentclass[addpoints]{exam}
\usepackage{amsmath}
\usepackage{amssymb}
\usepackage{color}
\usepackage[version=4]{mhchem}
\usepackage{siunitx}

% \printanswers

\newcommand{\event}{Matter, Bonding \& Periodic Trends}
\newcommand{\tournament}{}
\def\type{0}

\setlength\fillinlinelength{\textwidth}
\CorrectChoiceEmphasis{\color{red}\bfseries}
\SolutionEmphasis{\color{red}\bfseries}

\setlength\answerclearance{2pt}
\pagestyle{headandfoot}
\runningheadrule
\runningheader{\event}{\tournament}{\ifnum\type=1 Class Set Test \else Full Name:\hspace{1cm} \fi}
\runningfootrule
\runningfooter{}{Page \thepage\ of \numpages}{}

\begin{document}
\begin{center}
    {\huge Grade 11 Chemistry \\ \event
    \ifprintanswers \space Key
    \else \space \ifnum\type<2 Test \fi \ifnum\type=1 \textbf{(CLASS SET)} \fi
          \ifnum\type=2 Answer Sheet \fi
    \fi \\ \vspace{3mm}\tournament\vspace{1mm}}
\end{center}

\vspace{.25\textheight}

\begin{center}
    First Name: \ifprintanswers
    \underline{\hspace{3.5cm}}\textbf{KEY}\underline{\hspace{5.5cm}}\\\vspace{5mm}
    \else \underline{\hspace{10cm}}\\ \vspace{5mm} \fi
    Last Name: \underline{\hspace{10cm}}\\ \vspace{5mm}
\end{center}

\noindent\textbf{\underline{Directions:}}
\begin{itemize}
    \ifnum\type=1 \item \textbf{This test is a class set.} Do not write on this paper. \fi
    \item Show all work for full marks. A periodic table is provided.
    \item The test is designed to be completed in 75 minutes.
\end{itemize}
\begin{center}
\textit{For grading use only}\\
\multirowgradetable{1}[pages]
\end{center}

\newpage
\begin{questions}

\section*{Part A --- Multiple Choice (15 marks)}
\vspace{5mm}

\question[1] An atom of \ce{^{37}Cl} has:
\begin{choices}
    \choice 17 protons, 17 neutrons, 17 electrons
    \correctchoice 17 protons, 20 neutrons, 17 electrons
    \choice 17 protons, 20 neutrons, 20 electrons
    \choice 20 protons, 17 neutrons, 17 electrons
\end{choices}

\question[1] Which of the following correctly lists atomic radius in \textbf{decreasing} order?
\begin{choices}
    \choice $\text{Li} < \text{Na} < \text{K}$
    \correctchoice $\text{K} > \text{Na} > \text{Li}$
    \choice $\text{F} > \text{Cl} > \text{Br}$
    \choice $\text{Na} > \text{K} > \text{Rb}$
\end{choices}

\question[1] First ionization energy generally \textbf{increases} as you move:
\begin{choices}
    \choice Down a group
    \choice Across a period from right to left
    \correctchoice Across a period from left to right
    \choice Down and to the right on the periodic table
\end{choices}

\question[1] The most electronegative element on the periodic table is:
\begin{choices}
    \choice Oxygen
    \choice Chlorine
    \correctchoice Fluorine
    \choice Nitrogen
\end{choices}

\question[1] Which bond is most polar?
\begin{choices}
    \choice \ce{C-H}
    \choice \ce{N-H}
    \correctchoice \ce{O-H}
    \choice \ce{C-C}
\end{choices}

\question[1] The formula for the ionic compound formed between calcium and phosphate is:
\begin{choices}
    \choice \ce{CaPO4}
    \choice \ce{Ca2(PO4)3}
    \correctchoice \ce{Ca3(PO4)2}
    \choice \ce{Ca(PO4)2}
\end{choices}

\question[1] How many bonding pairs and lone pairs are in a molecule of \ce{NH3}?
\begin{choices}
    \choice 2 bonding pairs, 2 lone pairs
    \correctchoice 3 bonding pairs, 1 lone pair
    \choice 3 bonding pairs, 0 lone pairs
    \choice 4 bonding pairs, 0 lone pairs
\end{choices}

\question[1] According to VSEPR theory, the molecular geometry of \ce{H2O} is:
\begin{choices}
    \choice Linear
    \choice Trigonal planar
    \correctchoice Bent (V-shaped)
    \choice Tetrahedral
\end{choices}

\question[1] The molecular geometry of \ce{BF3} is:
\begin{choices}
    \choice Bent
    \correctchoice Trigonal planar
    \choice Trigonal pyramidal
    \choice T-shaped
\end{choices}

\question[1] Which of the following molecules is \textbf{nonpolar} despite having polar bonds?
\begin{choices}
    \choice \ce{HCl}
    \choice \ce{H2O}
    \correctchoice \ce{CO2}
    \choice \ce{NH3}
\end{choices}

\question[1] The strongest intermolecular force in liquid \ce{HF} is:
\begin{choices}
    \choice London dispersion forces
    \choice Dipole-dipole forces
    \correctchoice Hydrogen bonding
    \choice Ionic bonding
\end{choices}

\question[1] London dispersion forces are present in:
\begin{choices}
    \choice Only polar molecules
    \choice Only nonpolar molecules
    \correctchoice All molecules
    \choice Only ionic compounds
\end{choices}

\question[1] The ion \ce{S^{2-}} has:
\begin{choices}
    \choice 16 protons and 14 electrons
    \correctchoice 16 protons and 18 electrons
    \choice 14 protons and 16 electrons
    \choice 18 protons and 16 electrons
\end{choices}

\question[1] Which of the following has the highest melting point?
\begin{choices}
    \choice \ce{CH4}
    \choice \ce{HCl}
    \correctchoice \ce{NaCl}
    \choice \ce{H2O}
\end{choices}

\question[1] A double bond between two atoms consists of:
\begin{choices}
    \choice Two sigma bonds
    \correctchoice One sigma bond and one pi bond
    \choice Two pi bonds
    \choice One sigma bond and two pi bonds
\end{choices}


\section*{Part B --- Short Answer (33 marks)}

\question Answer the following questions about atomic structure and periodic trends.
\begin{parts}
    \part[2] Write the full electron configuration for sulfur (\ce{S}, $Z=16$) and for
    the sulfide ion (\ce{S^{2-}}).
    \part[3] Arrange the following in order of \textbf{increasing} first ionization energy
    and explain the trend: \ce{Na}, \ce{Mg}, \ce{Al}, \ce{Si}, \ce{P}.
    \part[3] Explain why the second ionization energy of sodium is much greater than its
    first ionization energy.
\end{parts}
\begin{solution}[9cm]
\textbf{(a)} \ce{S}: $1s^2\,2s^2\,2p^6\,3s^2\,3p^4$\quad
\ce{S^{2-}}: $1s^2\,2s^2\,2p^6\,3s^2\,3p^6$ (gains 2 electrons, isoelectronic with Ar)

\textbf{(b)} Increasing IE$_1$: \ce{Na} $<$ \ce{Mg} $<$ \ce{Al} $<$ \ce{Si} $<$ \ce{P}
(note: Mg has slightly higher IE than Al because Al's 3p electron is easier to remove than
Mg's filled 3s). General trend: IE increases across a period as nuclear charge increases
and atomic radius decreases.

\textbf{(c)} Sodium's first IE removes the single 3s electron, which is far from the nucleus
and well-shielded. The second IE removes a 2p electron from a much more stable, fully-filled
shell (noble gas configuration), requiring far more energy.
\end{solution}

\question Draw the Lewis structure for each of the following, then predict the molecular
geometry using VSEPR theory.
\begin{parts}
    \part[3] \ce{PCl3}
    \part[3] \ce{CO2}
    \part[2] \ce{CH4}
\end{parts}
\begin{solution}[9cm]
\textbf{(a)} \ce{PCl3}: P is the central atom with 3 bonding pairs to Cl and 1 lone pair on P.
Electron geometry: tetrahedral. Molecular geometry: \textbf{trigonal pyramidal}. Bond angles
$\approx 107^\circ$.

\textbf{(b)} \ce{CO2}: C forms two double bonds with O (total 4 bonding pairs, 0 lone pairs
on C). Electron geometry: linear. Molecular geometry: \textbf{linear}. Bond angle $180^\circ$.
The molecule is nonpolar because the two C=O dipoles cancel.

\textbf{(c)} \ce{CH4}: C forms 4 single bonds with H and has no lone pairs. Electron geometry:
tetrahedral. Molecular geometry: \textbf{tetrahedral}. Bond angles $\approx 109.5^\circ$.
\end{solution}

\question Ionic compounds and bonding.
\begin{parts}
    \part[2] Write the formula and name for the ionic compound formed between aluminum
    (\ce{Al^{3+}}) and sulfate (\ce{SO4^{2-}}).
    \part[3] Explain why ionic compounds have high melting points but dissolve in water to
    produce conducting solutions.
    \part[3] Compare the boiling points of \ce{H2O}, \ce{H2S}, and \ce{H2Se}.
    Explain the trend and account for the anomalously high boiling point of \ce{H2O}.
\end{parts}
\begin{solution}[9cm]
\textbf{(a)} \ce{Al2(SO4)3} — aluminum sulfate. (2 Al$^{3+}$ balanced by 3 SO$_4^{2-}$)

\textbf{(b)} Ionic compounds are held together by strong electrostatic attractions between
oppositely charged ions in a crystal lattice, requiring high temperatures to break these
forces (high melting point). In water, polar water molecules surround and stabilize individual
ions (solvation), breaking apart the lattice. The resulting ions are free to move and carry
charge (conducting solution).

\textbf{(c)} For the group 16 hydrides, boiling point generally increases with molar mass:
\ce{H2S} $< $ \ce{H2Se} due to increasing London dispersion forces. However, \ce{H2O} has a
far higher boiling point ($\SI{100}{\celsius}$) than expected because \ce{H2O} molecules form
strong \textbf{hydrogen bonds} (due to the highly electronegative O and O--H bonds), which
require much more energy to overcome.
\end{solution}

\question Polarity and intermolecular forces.
\begin{parts}
    \part[2] Identify the type(s) of intermolecular forces present in each: (i) \ce{Br2};\;
    (ii) \ce{HCl};\; (iii) \ce{CH3OH}.
    \part[3] Explain why ethanol (\ce{CH3CH2OH}) is miscible with water but hexane
    (\ce{C6H14}) is not, using the principle ``like dissolves like.''
    \part[3] Arrange the following in order of increasing boiling point and justify:
    \ce{Ne}, \ce{Ar}, \ce{Kr}, \ce{Xe}.
\end{parts}
\begin{solution}[9cm]
\textbf{(a)}
(i) \ce{Br2}: London dispersion forces only (nonpolar molecule).
(ii) \ce{HCl}: London dispersion + dipole-dipole forces (polar molecule).
(iii) \ce{CH3OH}: London dispersion + dipole-dipole + \textbf{hydrogen bonding}
(O--H group present).

\textbf{(b)} Ethanol is polar and can form hydrogen bonds with water (both have O--H groups),
so they are miscible. Hexane is a nonpolar hydrocarbon with only London dispersion forces.
Since water is highly polar, the strong H-bonds between water molecules are not replaced by
favorable interactions with nonpolar hexane, so hexane does not dissolve.

\textbf{(c)} Increasing boiling point: \ce{Ne} $<$ \ce{Ar} $<$ \ce{Kr} $<$ \ce{Xe}.
All noble gases have only London dispersion forces. As molar mass increases down the group,
more electrons mean greater polarizability, stronger temporary dipoles, and stronger LDF,
requiring more energy to vaporize.
\end{solution}

\end{questions}
\end{document}
